\chapter{Regulators, Deference and Cost Benefit Analysis}
\label{ch:regulation}

\section{The administrative state as a decision maker}

Much of modern law is made and applied not by courts but by agencies: environmental protection agencies setting emission limits, securities regulators writing disclosure rules, financial conduct authorities enforcing standards of conduct, health regulators approving drugs. Agencies make decisions under uncertainty about the effects of their rules, and courts review those decisions. The questions this book has asked of courts apply equally to agencies: what are the objectives, what are the probabilities, what are the losses, and is the decision coherent? This chapter examines cost benefit analysis as the agency's version of expected utility, the judicial review of agency decisions, and the Supreme Court's 2024 decision ending the doctrine of deference to agency interpretations of statutes, which can be read as a change in the weight given to one source of evidence about the meaning of law.\index{administrative law}\index{regulation}

\section{Cost benefit analysis}

\subsection{Expected value in government}

Since 1981, American executive orders (Executive Order 12291, then Executive Order 12866) have required executive agencies (Cabinet departments and independent agencies; independent regulatory agencies such as the SEC were excluded until Executive Order 14215 of 2025 extended review to them) to assess the costs and benefits of significant regulations, and since 1993 Executive Order 12866 has directed agencies to propose or adopt a regulation only upon a reasoned determination that its benefits justify its costs. Cost benefit analysis is expected value calculation (expected utility with a linear utility, that is, risk neutrality, which the precaution discussion below relaxes): the agency estimates the probability distribution of each effect, monetises it, and compares the expected net benefits of alternatives. The value of a statistical life\index{value of a statistical life}, discussed in Chapter~\ref{ch:becker}, converts reductions in mortality risk into money. Discounting converts future effects into present values.\index{cost benefit analysis}

\subsection{Uncertainty}

Many regulatory effects are deeply uncertain: the health effects of low doses of a chemical, the systemic risk reduced by a capital requirement, the climate damages avoided by an emissions limit. Cost benefit analysis under deep uncertainty requires distributions, not point estimates (subjective or elicited distributions where none is agreed), and the choice of distribution can drive the result. The Office of Management and Budget's Circular A-4 asks agencies to characterise uncertainty, and requires a formal quantitative analysis of it for rules with annual effects of one billion dollars or more \citep{omb2003a4}. (The 2023 revision of the circular was rescinded in 2025, which reinstated the 2003 version.) A Bayesian analysis reports the posterior distribution of net benefits and the probability that a rule's benefits exceed its costs, and it can incorporate expert elicitation where data are thin, in the manner of Smith's decision analysis \citep{smith2010}.

\subsection{The precautionary principle}

European law has given a prominent place to the precautionary principle, under which uncertainty about serious harm is not a reason to postpone measures to prevent it. Critics argue that the principle, taken literally, is incoherent, since both action and inaction carry risks, and that it can justify any regulation \citep{sunstein2005}. A decision theoretic reading makes it coherent: when the loss from an irreversible harm is very large relative to the cost of precaution, the expected loss from inaction can exceed that of action even at a modest probability of harm, which is the Hand formula of Chapter~\ref{ch:tort} at the scale of a society. Read this way, precaution is not an alternative to cost benefit analysis but a statement about the shape of the loss function, in particular about the weight on catastrophic and irreversible outcomes.

\section{Judicial review of costs and benefits}

Courts review agency decisions for reasonableness. In \emph{Business Roundtable v SEC} the District of Columbia Circuit in 2011 vacated a Securities and Exchange Commission rule giving shareholders access to company proxy materials to nominate directors, holding that the Commission had acted arbitrarily by failing adequately to assess the rule's economic effects \citep{businessroundtable2011}. In \emph{Michigan v EPA} the Supreme Court held in 2015 that the Environmental Protection Agency had acted unreasonably in deciding to regulate power plant emissions without considering cost, since reasonable regulation ordinarily requires "paying attention to the advantages and the disadvantages of agency decisions" \citep{michiganepa2015}.

These decisions require agencies to attend to costs and benefits, though \emph{Michigan v EPA} expressly declined to require a formal analysis in which every effect is monetised, and they put courts in the position of assessing the quality of that reasoning. A Bayesian standard of review would ask whether the agency identified the relevant effects, characterised their uncertainty, used defensible probabilities and values, and chose the alternative with the highest expected net benefit, or explained why it did not. This is a standard of coherence, not of correctness: the court does not substitute its own probabilities for the agency's, but asks whether the agency's decision follows from its stated probabilities and values.

\section{Deference after Loper Bright}

\subsection{Chevron and its end}

For forty years, under \emph{Chevron v Natural Resources Defense Council} \citep{chevron1984}, American courts deferred to an agency's reasonable interpretation of an ambiguous statute that the agency administered. In \emph{Loper Bright Enterprises v Raimondo} the Supreme Court in 2024 overruled \emph{Chevron}, holding that the Administrative Procedure Act requires courts to exercise their independent judgment in deciding whether an agency has acted within its statutory authority, and that courts may not defer to an agency interpretation merely because a statute is ambiguous \citep{loperbright2024}. Courts may still give weight to an agency's views, under the older doctrine of \emph{Skidmore v Swift}, according to the thoroughness evident in its consideration, the validity of its reasoning, its consistency and other factors that give it power to persuade \citep{skidmore1944}.\index{Loper Bright}\index{Chevron}

\subsection{A Bayesian reading}

Interpretation of a statute, like interpretation of a contract in Chapter~\ref{ch:contract}, is inference about meaning from text and context. Let $m$ range over possible meanings. The court forms a posterior $\Prob(m\given\text{text},\text{context})$ and adopts the meaning that minimises expected loss, often the most probable one, or one with high posterior probability. An agency's interpretation is evidence about meaning: the agency has expertise in the subject, may have participated in drafting, and has experience administering the statute. How much weight should that evidence receive?

Under \emph{Chevron}, the answer was, in effect, that the agency's view decided the question whenever the court's own posterior was not concentrated on one meaning: if the statute was ambiguous, the agency's reasonable reading prevailed. That is a rule of decision, not of belief: whenever the court was uncertain, it adopted the agency's reasonable reading whatever its own posterior. Under \emph{Skidmore}, as reaffirmed in \emph{Loper Bright}, the agency's interpretation is a signal with a likelihood ratio that depends on its quality: thorough, well reasoned and consistent interpretations carry more weight; hasty, poorly reasoned or shifting ones carry less. That is ordinary Bayesian updating on expert testimony, with the weight of the testimony depending on the expert's reliability, as for any expert in Chapter~\ref{ch:evidence}.

The Bayesian reading clarifies what changed. \emph{Loper Bright} did not make agency views irrelevant; it changed their status from a rule of decision to a source of evidence. The practical consequences depend on how much weight courts give agency interpretations as evidence, which will vary with the quality of the agency's reasoning. Agencies that explain their interpretations thoroughly, with attention to text, purpose and consequences, will continue to persuade; those that rely on authority are less likely to. The decision also increases uncertainty about the law, since meanings once settled by deference may now be relitigated, and that uncertainty has costs for regulated parties that the theory of Chapter~\ref{ch:systemcosts} would count.

\section{Enforcement as a decision}

Agencies also enforce. A securities or conduct regulator decides which firms to investigate, which cases to bring, and whether to settle. Each decision is the subject of the economics of Chapter~\ref{ch:becker}: the probability of detection and the sanction together determine deterrence, and enforcement resources are limited. In the United Kingdom, the Financial Conduct Authority's Consumer Duty, in force from 2023, requires firms to act to deliver good outcomes for retail customers, a standard expressed in terms of outcomes rather than processes, which invites evidence about outcomes, such as complaint rates, value for money and customer understanding, and therefore invites statistical monitoring. A regulator monitoring outcomes faces the screening problem of the cartel screens in Chapter~\ref{ch:antitrust}: a flag has a likelihood ratio, the posterior probability of misconduct it yields depends on the base rate, and most flags in a low base rate population will be false unless the false positive rate is very small.

\section{Regulating AI in law}

The regulation of artificial intelligence, to which Chapters~\ref{ch:lindleyai} and~\ref{ch:xai} return, is itself an example of regulatory decision making under deep uncertainty. The European Union's AI Act \citep{euaiact2024} classifies systems by risk and imposes obligations proportional to it; other jurisdictions have chosen lighter or sectoral approaches. The choice depends on probabilities that are poorly known, about the harms and benefits of systems that do not yet exist, and on values about innovation, safety and rights. Cost benefit analysis of AI regulation is hard, but the framework still applies: state the objectives, characterise the uncertainty, compare the alternatives, and revise as evidence accumulates. Rules that build in review and revision, adapting as posteriors change, are preferable under such uncertainty to rules that assume the answers are known.

\section*{Notes and Questions}
\addcontentsline{toc}{section}{Notes and Questions}
\begin{enumerate}[label=\arabic*.]
\item A regulation costs one hundred million dollars a year and reduces annual mortality risk by one in a million for a population of thirty million. At a value of a statistical life of ten million, do its benefits exceed its costs? By how much would the risk reduction have to be overestimated for the answer to change?
\item Formalise the precautionary principle as a statement about the loss function. Under what conditions does it imply action at a low probability of harm?
\item Compare \emph{Chevron} deference and \emph{Skidmore} weight. Which is a rule for updating a court's posterior on the meaning of a statute, and which is a rule of decision? What difference does that make when the court's posterior is spread over two meanings?
\item Read \emph{Business Roundtable v SEC}. What would a coherent cost benefit analysis of the proxy access rule have required?
\item A conduct regulator's outcome monitoring flags five per cent of firms; misconduct occurs at one per cent of firms; the monitor catches eighty per cent of misconduct. What fraction of flagged firms have engaged in misconduct?
\end{enumerate}

\section*{Further reading}
\emph{Loper Bright} \citep{loperbright2024}; \emph{Skidmore} \citep{skidmore1944}; \emph{Michigan v EPA} \citep{michiganepa2015}; \emph{Business Roundtable v SEC} \citep{businessroundtable2011}; Smith on decision analysis \citep{smith2010}.
